\section{RQ2: Corrections and Improvements of POEMS}
\label{sec:rq2}

\subsection{Effect of the Corrections}

Table~\ref{tab:rq2fixes} compares POEMS before and after the corrections of Section~\ref{sec:defects}, both with the fast immutable implementation. The corrections change the behaviour of the algorithm in several ways at once, and their combined effect within 30 seconds is mixed. The dead space is significantly smaller on n300 and significantly larger on n100; on the remaining benchmarks the differences are not significant (smaller means on xerox, hp and ami49, larger on ami33 and n200).

\begin{table}
\centering
\small
\begin{tabular}{@{}lcc@{}}
\toprule
Benchmark & Before corrections & After corrections \\
\midrule
apte & $0.78_{\pm 0.00}\ (0.78)$ & $0.78_{\pm 0.00}\ (0.78)$ \\
xerox & $3.07_{\pm 0.41}\ (2.34)$ & $2.79_{\pm 0.37}\ (2.34)$ \\
hp & $4.43_{\pm 1.08}\ (2.28)$ & $3.79_{\pm 1.76}\ (1.32)$ \\
ami33 & $2.58_{\pm 0.46}\ (1.73)$ & $2.98_{\pm 0.59}\ (2.22)$ \\
ami49 & $3.37_{\pm 0.64}\ (2.31)$ & $2.95_{\pm 0.39}\ (2.34)$ \\
n100 & $4.15_{\pm 0.60}\ (3.49)$ & $5.00_{\pm 0.60}\ (4.43)^{\blacktriangle}$ \\
n200 & $7.30_{\pm 0.47}\ (6.67)$ & $7.66_{\pm 0.42}\ (6.73)$ \\
n300 & $9.00_{\pm 0.51}\ (7.99)$ & $\mathbf{8.12_{\pm 0.56}\ (7.14)}$ \\
\bottomrule
\end{tabular}
\caption{Dead space [\%] before and after the corrections of Section~\ref{sec:defects} (30\,s, 8 seeds)}
\label{tab:rq2fixes}
\end{table}


Two of the corrections change the search itself. The corrected prototype is a different, deterministic starting point of all runs. The corrected identity check makes steps that only rotate modules possible. To separate them on n100, each was reverted alone. With the original prototype, the mean dead space on n100 is 4.49\,\% (8 seeds); with the original identity check, it is 4.36\,\%. Both lie between the results before (4.15\,\%) and after (5.00\,\%) the corrections, so both contribute to the difference, and neither difference is significant on its own.

Table~\ref{tab:rotation} shows how the identity check changes the iterations. With the original check, a step that only rotates modules is rejected, so the genetic algorithm has to find a structural change. When it does not find a good one, the iteration ends with a worse step, which is not applied (9--16\,\% of the iterations on ami33, n100 and n300). With the corrected check, a rotation of a module that does not touch the boundary of the floorplan is a step of equal fitness, which is found almost immediately (on ami33, n100 and n300, the mean generation of the best individual drops from 128--175 to 23--32) and accepted. The share of improving iterations does not change. In other words, the correction turns rejected iterations into moves along plateaus. Whether such a move helps depends on the instance, which agrees with the mixed results of Table~\ref{tab:rq2fixes}.

\begin{table}
\centering
\small
\begin{tabular}{@{}lcccc@{}}
\toprule
 & \multicolumn{2}{c}{iterations improving / equal / worse [\%]} & \multicolumn{2}{c}{generation of the best (mean)} \\
Benchmark & rotation ignored & corrected & rotation ignored & corrected \\
\midrule
xerox & 0.4 / 97.8 / 1.8 & 0.4 / 98.0 / 1.5 & 50 & 40 \\
ami33 & 0.9 / 87.5 / 11.6 & 0.6 / 99.0 / 0.4 & 138 & 23 \\
n100 & 3.0 / 88.0 / 9.0 & 2.9 / 96.9 / 0.2 & 128 & 24 \\
n300 & 5.3 / 78.9 / 15.9 & 4.9 / 94.9 / 0.2 & 175 & 32 \\
\bottomrule
\end{tabular}
\caption{Outcome of the iterations of POEMS with the original identity check, which ignores the orientation of the modules, and with the corrected one (20\,s, seed 1)}
\label{tab:rotation}
\end{table}


The corrections are kept in all further experiments: they make the implementation agree with the algorithm described in Chapter~\ref{sec:approach}, and the remaining defects (the hang for $N = 1$, the unreachable last node, the first iteration that could make the prototype worse) are plain errors.

\subsection{Improvements of the Genetic Algorithm}

All variants in this subsection use the corrections and the array evaluation, which gives the same results as the immutable evaluation and makes POEMS another 1.7 to 3.7 times faster (Table~\ref{tab:speed}). Each variant changes one setting of the reference configuration, which uses the parameters of the original edition ($G = 1000$, $N = 3$, $S = 50$). Table~\ref{tab:rq2ga} shows the results. The rule for the {\em improved POEMS} of the final comparison was fixed before the experiment: it combines all modifications whose score is better than that of the reference (of the reduced numbers of generations, only the better one).

\begin{table}
\centering
\footnotesize
\setlength{\tabcolsep}{4pt}
\begin{tabular}{@{}lcccccccc@{}}
\toprule
Benchmark & reference & $G = 300$ & $G = 100$ & no dupl. & move & patience & first & LA 20 \\
\midrule
xerox & $2.66_{\pm 0.35}$ & $2.61_{\pm 0.38}$ & $2.74_{\pm 0.34}$ & $2.92_{\pm 0.37}$ & $3.01_{\pm 0.32}$ & $2.93_{\pm 0.43}$ & $2.94_{\pm 0.38}$ & $2.70_{\pm 0.34}$ \\
hp & $2.41_{\pm 1.59}$ & $4.23_{\pm 0.78}$ & $2.14_{\pm 1.17}$ & $3.23_{\pm 1.69}$ & $2.64_{\pm 1.87}$ & $2.83_{\pm 1.59}$ & $2.41_{\pm 1.59}$ & $3.81_{\pm 1.35}$ \\
ami33 & $2.74_{\pm 0.47}$ & $\mathbf{1.89_{\pm 0.48}}$ & $\mathbf{1.92_{\pm 0.18}}$ & $2.58_{\pm 0.29}$ & $2.27_{\pm 0.45}$ & $\mathbf{1.83_{\pm 0.37}}$ & $2.50_{\pm 0.16}$ & $2.74_{\pm 0.47}$ \\
ami49 & $2.33_{\pm 0.23}$ & $2.40_{\pm 0.28}$ & $2.45_{\pm 0.35}$ & $2.90_{\pm 0.48}$ & $2.82_{\pm 0.65}$ & $2.47_{\pm 0.23}$ & $2.71_{\pm 0.29}^{\blacktriangle}$ & $2.55_{\pm 0.28}$ \\
n100 & $3.72_{\pm 0.45}$ & $\mathbf{3.09_{\pm 0.28}}$ & $\mathbf{2.97_{\pm 0.25}}$ & $3.53_{\pm 0.62}$ & $3.79_{\pm 0.79}$ & $3.16_{\pm 0.25}$ & $3.73_{\pm 0.39}$ & $3.83_{\pm 0.58}$ \\
n200 & $5.60_{\pm 0.51}$ & $\mathbf{4.99_{\pm 0.24}}$ & $\mathbf{4.47_{\pm 0.11}}$ & $5.57_{\pm 0.49}$ & $6.20_{\pm 0.56}$ & $\mathbf{4.33_{\pm 0.42}}$ & $5.89_{\pm 0.34}$ & $6.18_{\pm 0.42}$ \\
n300 & $6.51_{\pm 0.53}$ & $5.86_{\pm 0.58}$ & $\mathbf{5.57_{\pm 0.40}}$ & $6.19_{\pm 0.57}$ & $6.29_{\pm 0.40}$ & $\mathbf{5.42_{\pm 0.43}}$ & $6.67_{\pm 0.32}$ & $6.65_{\pm 0.29}$ \\
\midrule
score & 1.188 & 1.154 & 1.030 & 1.265 & 1.240 & 1.077 & 1.229 & 1.315 \\
\bottomrule
\end{tabular}
\caption{Modifications of the genetic algorithm of POEMS, each applied alone to the reference configuration (array evaluation, $G = 1000$, $N = 3$, $S = 50$): fewer generations per iteration, discarding duplicate children, the move node action, ending an iteration after 100 generations without improvement or at the first improvement of the prototype, and late acceptance of prototypes with a history of 20 iterations (dead space [\%], 30\,s, seeds 1--6)}
\label{tab:rq2ga}
\end{table}


The number of generations per iteration matters most. With $G = 100$ instead of 1000, an iteration costs 320 instead of 1856 evaluations, POEMS makes 5.6 times more iterations in the same time, and its dead space is significantly smaller on ami33, n100, n200 and n300. This agrees with the analysis of Section~\ref{sec:gaanalysis}: the best step of an iteration is usually found early, so the later generations mostly delay the next iteration. $G = 300$ lies in between. Discarding duplicate children saves 42\,\% of the evaluations of an iteration and gives up to 45\,\% more iterations on the large instances, where it is slightly better, but it is worse on xerox, hp and ami49; none of the differences is significant. A plausible explanation is that duplicates of good sequences are not only waste: they increase the share of the good sequences in the population, so discarding them weakens the selection. The move node action does not help either. Ending an iteration after 100 generations without improvement (patience) has an effect similar to $G = 100$ (about 330 evaluations per iteration) and is significantly better on ami33, n200 and n300. Ending it at the first improvement of the prototype is worse: because few iterations improve the prototype, the iterations still run almost all 1000 generations (about 1840 instead of 1856 evaluations per iteration), and when one improves it, the improving step is applied before a better one can be found. Late acceptance of prototypes is worse as well. It rarely has an effect, because almost all iterations end with a step of equal fitness, so the history mostly contains the current fitness; on ami33 it did not change the result of any run.

By the rule stated above, the improved POEMS combines $G = 100$ with the patience of 100 generations. With $G = 100$, however, the patience never ends an iteration early, so the improved POEMS is simply POEMS with the corrections, the array evaluation and $G = 100$, and its screening results are those of the column $G = 100$ of Table~\ref{tab:rq2ga}.

\paragraph{Answer to RQ2.} The genetic algorithm of POEMS uses its evaluations inefficiently. Almost half of its children duplicate individuals of the population, the best step of an iteration is usually present in the initial random population, and only a few percent of the iterations improve the prototype. The corrections of the defects change the behaviour of the algorithm, mainly by allowing pure rotations, but their combined effect on the results is mixed. Of the evaluated modifications, only those that shorten the iterations help: with $G = 100$, POEMS is significantly better on four of the seven benchmarks, and with the patience of 100 generations on three. Duplicate elimination, the move node action, the first-improvement rule and late acceptance of prototypes do not improve it.
